\section{Related Work}
\label{sec:related_work}
Research on communication spans many fields, e.g. linguistics, psychology, evolution and AI.  In AI, it is split along a few axes: a)~predefined or learned communication protocols, b)~planning or learning methods, c)~evolution or RL, and d)~cooperative or competitive settings. 

Given the topic of our paper, we focus on related work that deals with the cooperative learning of communication protocols. Out of the plethora of work on multi-agent RL with communication, e.g., \cite{tan1993multi,melo2011querypomdp,Panait2005387,Zhang20131101,kasai2008learning}, only a few fall into this category. Most assume a pre-defined communication protocol, rather than trying to learn protocols. 
One exception is the work of \citet{kasai2008learning}, in which tabular Q-learning agents have to learn the content of a message to solve a predator-prey task with communication.
Another example of open-ended communication learning in a multi-agent task is given in \cite{giles2002learning}. Here evolutionary methods are used for learning the protocols which are evaluated on a similar predator-prey task. Their approach uses a fitness function that is carefully designed to accelerate learning.
In general, heuristics and  handcrafted rules have prevailed widely in this line of research. Moreover, typical tasks have been necessarily small so that global optimisation methods, such as evolutionary algorithms, can be applied. The use of deep representations and gradient-based optimisation as advocated in this paper is an important departure, essential for scalability and further progress. A similar rationale is provided in~\cite{gregor2015draw}, another example of making an RL problem end-to-end differentiable.

Finally, we consider discrete communication channels. One of the key components of our methods is the signal binarisation during the decentralised execution. This is related to recent research on fitting neural networks in low-powered devices with memory and computational limitations using binary weights, e.g. \cite{courbariaux2016binarynet}, and previous works on discovering binary codes for documents~\cite{hinton2011discovering}.